% =====================================================================
% ΠΡΩΤΟ ΔΙΚΟ ΣΟΥ ΚΕΦΑΛΑΙΟ — ΑΣΚΗΣΗ ΧΡΗΣΗΣ ΤΟΥ TEMPLATE
% =====================================================================
% Το αρχείο ΔΕΝ βρίσκεται ακόμη στο book-content.tex και ΔΕΝ τυπώνεται.
% Για να το εμφανίσεις, γράψε εκεί: % =====================================================================
% ΠΡΩΤΟ ΔΙΚΟ ΣΟΥ ΚΕΦΑΛΑΙΟ — ΑΣΚΗΣΗ ΧΡΗΣΗΣ ΤΟΥ TEMPLATE
% =====================================================================
% Το αρχείο ΔΕΝ βρίσκεται ακόμη στο book-content.tex και ΔΕΝ τυπώνεται.
% Για να το εμφανίσεις, γράψε εκεί: % =====================================================================
% ΠΡΩΤΟ ΔΙΚΟ ΣΟΥ ΚΕΦΑΛΑΙΟ — ΑΣΚΗΣΗ ΧΡΗΣΗΣ ΤΟΥ TEMPLATE
% =====================================================================
% Το αρχείο ΔΕΝ βρίσκεται ακόμη στο book-content.tex και ΔΕΝ τυπώνεται.
% Για να το εμφανίσεις, γράψε εκεί: % =====================================================================
% ΠΡΩΤΟ ΔΙΚΟ ΣΟΥ ΚΕΦΑΛΑΙΟ — ΑΣΚΗΣΗ ΧΡΗΣΗΣ ΤΟΥ TEMPLATE
% =====================================================================
% Το αρχείο ΔΕΝ βρίσκεται ακόμη στο book-content.tex και ΔΕΝ τυπώνεται.
% Για να το εμφανίσεις, γράψε εκεί: \include{chapters/askisi-dokimis}
% Δεν χρειάζεται να ορίσεις σε αυτό το αρχείο χρώματα ή περιθώρια.
% =====================================================================

% 1. ΑΡΧΗ ΚΕΦΑΛΑΙΟΥ: τίτλος, υπότιτλος, εικόνα, πηγή, διδακτικοί στόχοι.
\BookChapterPhoto
  {Το πρώτο μου κεφάλαιο}
  {Δοκιμή των βασικών εντολών}
  {assets/chapter-photo-sample.png}
  {Εικόνα προσωρινή — μόνο για τη δοκιμή}
  {
    \item Να γράφω κείμενο και μαθηματικούς τύπους.
    \item Να προσθέτω λυμένα παραδείγματα.
    \item Να προσθέτω άλυτες ασκήσεις.
    \item Να κάνω παραπομπές και να δημιουργώ PDF.
  }

% 2. ΚΥΡΙΑ ΕΝΟΤΗΤΑ: η μπλε λωρίδα εμφανίζεται αυτόματα.
\BookSection{Θεωρία}
% Η κανονική \section αποκτά αριθμό χωρίς χειροκίνητη πληκτρολόγηση.
\section{Ένας βασικός ορισμός}
Γράψε εδώ τον δικό σου ορισμό. Το γράμμα \(x\) βρίσκεται μέσα στο κείμενο.

% 3. ΕΞΙΣΩΣΗ: \[ ... \] σημαίνει μαθηματικός τύπος σε δική του γραμμή.
\[\Delta x=x_{\mathrm{τελ}}-x_{\mathrm{αρχ}}.\]

% 4. ΛΥΜΕΝΟ ΠΑΡΑΔΕΙΓΜΑ: κάθε \begin πρέπει να έχει το αντίστοιχο \end.
\BookSection{Λυμένα παραδείγματα}
\begin{workedexample}{Το πρώτο μου λυμένο παράδειγμα}
  \begin{ProblemBox}
  Ένα κινητό περνά από την αρχική θέση \(x_A=1\,\mathrm{m}\)
  στην τελική θέση \(x_B=4\,\mathrm{m}\). Να υπολογιστεί η μετατόπιση.
  \end{ProblemBox}

  \begin{ThinkBox}
  Υπολογίζω τη διαφορά τελικής και αρχικής θέσης.
  \end{ThinkBox}

  \SolutionHeading
  \[\Delta x=4-1=3\,\mathrm{m}.\]
\end{workedexample}

% 5. ΑΛΥΤΗ: προαιρετικό ID σε [ ], κενός ορατός τίτλος σε {}.
\BookSection{Άλυτες ασκήσεις}
\begin{exercise}[MY-FIRST-001]{}
Να υπολογίσετε τη μετατόπιση κινητού από τη θέση \(x_A=-2\,\mathrm{m}\)
στη θέση \(x_B=5\,\mathrm{m}\).
\end{exercise}

% 6. ΘΕΩΡΙΑ: η εντολή TrueFalseItem κρατά μαζί τα δύο τετραγωνάκια.
\BookSection{Ερωτήσεις θεωρίας}
\BookQuestionHeading{Σωστό / Λάθος}
\begin{enumerate}
  \TrueFalseItem{Η μετατόπιση μπορεί να είναι αρνητική.}
  \TrueFalseItem{Το διάστημα είναι πάντοτε ίσο με τη μετατόπιση.}
\end{enumerate}

% Δεν χρειάζεται να ορίσεις σε αυτό το αρχείο χρώματα ή περιθώρια.
% =====================================================================

% 1. ΑΡΧΗ ΚΕΦΑΛΑΙΟΥ: τίτλος, υπότιτλος, εικόνα, πηγή, διδακτικοί στόχοι.
\BookChapterPhoto
  {Το πρώτο μου κεφάλαιο}
  {Δοκιμή των βασικών εντολών}
  {assets/chapter-photo-sample.png}
  {Εικόνα προσωρινή — μόνο για τη δοκιμή}
  {
    \item Να γράφω κείμενο και μαθηματικούς τύπους.
    \item Να προσθέτω λυμένα παραδείγματα.
    \item Να προσθέτω άλυτες ασκήσεις.
    \item Να κάνω παραπομπές και να δημιουργώ PDF.
  }

% 2. ΚΥΡΙΑ ΕΝΟΤΗΤΑ: η μπλε λωρίδα εμφανίζεται αυτόματα.
\BookSection{Θεωρία}
% Η κανονική \section αποκτά αριθμό χωρίς χειροκίνητη πληκτρολόγηση.
\section{Ένας βασικός ορισμός}
Γράψε εδώ τον δικό σου ορισμό. Το γράμμα \(x\) βρίσκεται μέσα στο κείμενο.

% 3. ΕΞΙΣΩΣΗ: \[ ... \] σημαίνει μαθηματικός τύπος σε δική του γραμμή.
\[\Delta x=x_{\mathrm{τελ}}-x_{\mathrm{αρχ}}.\]

% 4. ΛΥΜΕΝΟ ΠΑΡΑΔΕΙΓΜΑ: κάθε \begin πρέπει να έχει το αντίστοιχο \end.
\BookSection{Λυμένα παραδείγματα}
\begin{workedexample}{Το πρώτο μου λυμένο παράδειγμα}
  \begin{ProblemBox}
  Ένα κινητό περνά από την αρχική θέση \(x_A=1\,\mathrm{m}\)
  στην τελική θέση \(x_B=4\,\mathrm{m}\). Να υπολογιστεί η μετατόπιση.
  \end{ProblemBox}

  \begin{ThinkBox}
  Υπολογίζω τη διαφορά τελικής και αρχικής θέσης.
  \end{ThinkBox}

  \SolutionHeading
  \[\Delta x=4-1=3\,\mathrm{m}.\]
\end{workedexample}

% 5. ΑΛΥΤΗ: προαιρετικό ID σε [ ], κενός ορατός τίτλος σε {}.
\BookSection{Άλυτες ασκήσεις}
\begin{exercise}[MY-FIRST-001]{}
Να υπολογίσετε τη μετατόπιση κινητού από τη θέση \(x_A=-2\,\mathrm{m}\)
στη θέση \(x_B=5\,\mathrm{m}\).
\end{exercise}

% 6. ΘΕΩΡΙΑ: η εντολή TrueFalseItem κρατά μαζί τα δύο τετραγωνάκια.
\BookSection{Ερωτήσεις θεωρίας}
\BookQuestionHeading{Σωστό / Λάθος}
\begin{enumerate}
  \TrueFalseItem{Η μετατόπιση μπορεί να είναι αρνητική.}
  \TrueFalseItem{Το διάστημα είναι πάντοτε ίσο με τη μετατόπιση.}
\end{enumerate}

% Δεν χρειάζεται να ορίσεις σε αυτό το αρχείο χρώματα ή περιθώρια.
% =====================================================================

% 1. ΑΡΧΗ ΚΕΦΑΛΑΙΟΥ: τίτλος, υπότιτλος, εικόνα, πηγή, διδακτικοί στόχοι.
\BookChapterPhoto
  {Το πρώτο μου κεφάλαιο}
  {Δοκιμή των βασικών εντολών}
  {assets/chapter-photo-sample.png}
  {Εικόνα προσωρινή — μόνο για τη δοκιμή}
  {
    \item Να γράφω κείμενο και μαθηματικούς τύπους.
    \item Να προσθέτω λυμένα παραδείγματα.
    \item Να προσθέτω άλυτες ασκήσεις.
    \item Να κάνω παραπομπές και να δημιουργώ PDF.
  }

% 2. ΚΥΡΙΑ ΕΝΟΤΗΤΑ: η μπλε λωρίδα εμφανίζεται αυτόματα.
\BookSection{Θεωρία}
% Η κανονική \section αποκτά αριθμό χωρίς χειροκίνητη πληκτρολόγηση.
\section{Ένας βασικός ορισμός}
Γράψε εδώ τον δικό σου ορισμό. Το γράμμα \(x\) βρίσκεται μέσα στο κείμενο.

% 3. ΕΞΙΣΩΣΗ: \[ ... \] σημαίνει μαθηματικός τύπος σε δική του γραμμή.
\[\Delta x=x_{\mathrm{τελ}}-x_{\mathrm{αρχ}}.\]

% 4. ΛΥΜΕΝΟ ΠΑΡΑΔΕΙΓΜΑ: κάθε \begin πρέπει να έχει το αντίστοιχο \end.
\BookSection{Λυμένα παραδείγματα}
\begin{workedexample}{Το πρώτο μου λυμένο παράδειγμα}
  \begin{ProblemBox}
  Ένα κινητό περνά από την αρχική θέση \(x_A=1\,\mathrm{m}\)
  στην τελική θέση \(x_B=4\,\mathrm{m}\). Να υπολογιστεί η μετατόπιση.
  \end{ProblemBox}

  \begin{ThinkBox}
  Υπολογίζω τη διαφορά τελικής και αρχικής θέσης.
  \end{ThinkBox}

  \SolutionHeading
  \[\Delta x=4-1=3\,\mathrm{m}.\]
\end{workedexample}

% 5. ΑΛΥΤΗ: προαιρετικό ID σε [ ], κενός ορατός τίτλος σε {}.
\BookSection{Άλυτες ασκήσεις}
\begin{exercise}[MY-FIRST-001]{}
Να υπολογίσετε τη μετατόπιση κινητού από τη θέση \(x_A=-2\,\mathrm{m}\)
στη θέση \(x_B=5\,\mathrm{m}\).
\end{exercise}

% 6. ΘΕΩΡΙΑ: η εντολή TrueFalseItem κρατά μαζί τα δύο τετραγωνάκια.
\BookSection{Ερωτήσεις θεωρίας}
\BookQuestionHeading{Σωστό / Λάθος}
\begin{enumerate}
  \TrueFalseItem{Η μετατόπιση μπορεί να είναι αρνητική.}
  \TrueFalseItem{Το διάστημα είναι πάντοτε ίσο με τη μετατόπιση.}
\end{enumerate}

% Δεν χρειάζεται να ορίσεις σε αυτό το αρχείο χρώματα ή περιθώρια.
% =====================================================================

% 1. ΑΡΧΗ ΚΕΦΑΛΑΙΟΥ: τίτλος, υπότιτλος, εικόνα, πηγή, διδακτικοί στόχοι.
\BookChapterPhoto
  {Το πρώτο μου κεφάλαιο}
  {Δοκιμή των βασικών εντολών}
  {assets/chapter-photo-sample.png}
  {Εικόνα προσωρινή — μόνο για τη δοκιμή}
  {
    \item Να γράφω κείμενο και μαθηματικούς τύπους.
    \item Να προσθέτω λυμένα παραδείγματα.
    \item Να προσθέτω άλυτες ασκήσεις.
    \item Να κάνω παραπομπές και να δημιουργώ PDF.
  }

% 2. ΚΥΡΙΑ ΕΝΟΤΗΤΑ: η μπλε λωρίδα εμφανίζεται αυτόματα.
\BookSection{Θεωρία}
% Η κανονική \section αποκτά αριθμό χωρίς χειροκίνητη πληκτρολόγηση.
\section{Ένας βασικός ορισμός}
Γράψε εδώ τον δικό σου ορισμό. Το γράμμα \(x\) βρίσκεται μέσα στο κείμενο.

% 3. ΕΞΙΣΩΣΗ: \[ ... \] σημαίνει μαθηματικός τύπος σε δική του γραμμή.
\[\Delta x=x_{\mathrm{τελ}}-x_{\mathrm{αρχ}}.\]

% 4. ΛΥΜΕΝΟ ΠΑΡΑΔΕΙΓΜΑ: κάθε \begin πρέπει να έχει το αντίστοιχο \end.
\BookSection{Λυμένα παραδείγματα}
\begin{workedexample}{Το πρώτο μου λυμένο παράδειγμα}
  \begin{ProblemBox}
  Ένα κινητό περνά από την αρχική θέση \(x_A=1\,\mathrm{m}\)
  στην τελική θέση \(x_B=4\,\mathrm{m}\). Να υπολογιστεί η μετατόπιση.
  \end{ProblemBox}

  \begin{ThinkBox}
  Υπολογίζω τη διαφορά τελικής και αρχικής θέσης.
  \end{ThinkBox}

  \SolutionHeading
  \[\Delta x=4-1=3\,\mathrm{m}.\]
\end{workedexample}

% 5. ΑΛΥΤΗ: προαιρετικό ID σε [ ], κενός ορατός τίτλος σε {}.
\BookSection{Άλυτες ασκήσεις}
\begin{exercise}[MY-FIRST-001]{}
Να υπολογίσετε τη μετατόπιση κινητού από τη θέση \(x_A=-2\,\mathrm{m}\)
στη θέση \(x_B=5\,\mathrm{m}\).
\end{exercise}

% 6. ΘΕΩΡΙΑ: η εντολή TrueFalseItem κρατά μαζί τα δύο τετραγωνάκια.
\BookSection{Ερωτήσεις θεωρίας}
\BookQuestionHeading{Σωστό / Λάθος}
\begin{enumerate}
  \TrueFalseItem{Η μετατόπιση μπορεί να είναι αρνητική.}
  \TrueFalseItem{Το διάστημα είναι πάντοτε ίσο με τη μετατόπιση.}
\end{enumerate}
